% Generated draft: requires mathematical and source review.
\documentclass[11pt]{article}
\usepackage[margin=25mm]{geometry}
\usepackage{fontspec}
\setmainfont{Latin Modern Roman}
\usepackage{amsmath,amssymb,mathtools,unicode-math}
\setmathfont{Latin Modern Math}
\usepackage{enumitem,needspace,xurl,hyperref}
\setlength{\parindent}{0pt}
\setlength{\parskip}{5pt}
\setlength{\emergencystretch}{3em}
\newcommand{\N}{\mathbb{N}}
\newcommand{\Z}{\mathbb{Z}}
\newcommand{\Q}{\mathbb{Q}}
\newcommand{\R}{\mathbb{R}}
\newcommand{\C}{\mathbb{C}}
\newcommand{\F}{\mathbb{F}}
\newcommand{\A}{\mathbb{A}}
\newcommand{\PP}{\mathbb{P}}
\newcommand{\E}{\mathbb{E}}
\newcommand{\eps}{\varepsilon}
\newcommand{\dd}{\,\mathrm{d}}
\begin{document}
\section*{UnsolvedMath ID: 30004421}
\textbf{Source problem code:} catalogue entry
\subsection{Definitions and mathematical statement}
Let \(G\) be a classical group acting linearly on a vector space \(V\) over its ground field \(k\). The supplied record does not specify \(k\), the dimension of \(V\), the precise classical group, or the additional hypotheses intended by the phrase ``Aschbacher--Scott setting.'' These remain unresolved parameters of the problem.

A group \(S\) is quasisimple if \([S,S]=S\) and \(S/Z(S)\) is a nonabelian simple group, where \(Z(S)\) denotes its center. The action of \(S\) on \(V\) is irreducible over \(k\) if \(V\ne 0\) and its only \(S\)-invariant \(k\)-linear subspaces are \(0\) and \(V\). This does not impose absolute irreducibility.

For a proper quasisimple subgroup \(S<G\) acting irreducibly on \(V\), define its normalizer by
\[
M=N_G(S)=\{g\in G:gSg^{-1}=S\}.
\]
A subgroup \(M\) is maximal in \(G\) if \(M<G\) and every subgroup \(H\) satisfying \(M\le H\le G\) satisfies \(H=M\) or \(H=G\).

\textbf{Problem as supplied.} Determine necessary and sufficient conditions, within the intended Aschbacher--Scott setting, for \(N_G(S)\) to be a maximal subgroup of \(G\). Properness of \(S\) alone does not supply the required properness condition on \(N_G(S)\); that condition is part of the requested maximality.

The record describes this question as seeking a converse, but supplies neither the theorem to be reversed nor its hypotheses. Consequently, this is a characterization question with an unresolved scope, rather than a uniquely specified conjectural implication.
\subsection{Short English statement}
The question asks when the normalizer of an irreducibly acting quasisimple subgroup is maximal in its ambient classical group. The intended Aschbacher--Scott hypotheses, ground field, representation, and precise family of groups must be clarified before a fully specified characterization can be formulated.
\subsection{Sources}
\begin{itemize}
\item Local UnsolvedMath registry, permanent ID \(30004421\), source problem code \texttt{catalogue entry}, title ``Maximality of Normalizers of Quasisimple Subgroups,'' in the Oberwolfach Reports Open Problems collection. No individual proposer is identified; the proposed year is recorded as 2020.
\item Supplied source citation: \emph{Mini-Workshop: Kronecker, Plethysm, and Sylow Branching Coefficients and their Applications to Complexity Theory} (2021), Oberwolfach Reports, pp.\ 517--568, DOI \texttt{10.4171/owr/2020/9}. Supplied source URL and literature-source link labelled ``original OWR report'': \url{https://doi.org/10.4171/owr/2020/9}. The background describes these reports as documenting mathematical workshops and their open-problem discussions.
\item The supplied literature assessment was checked on \texttt{2026-08-22}. This is supplied dated curation, not independently verified here. It labels the problem open and states that the extract asks generally for a converse concerning maximal subgroups without specifying a precise theorem, so no uniquely identifiable resolved claim can be assessed. The supplied curation expressly calls for independent verification of time-sensitive status and cited claims.
\item The registry labels the statement \texttt{corrected\_verified} and reports a manual semantic, typographic, and LaTeX review claiming self-containment and preservation of the mathematical target. These are supplied editorial assertions; they do not resolve the missing mathematical hypotheses or constitute independent mathematical verification.
\end{itemize}
\end{document}
